\documentclass[11pt,a4paper]{article}
\usepackage{fontspec}
\usepackage{geometry}
\geometry{margin=2.5cm}
\usepackage{graphicx}
\usepackage{booktabs}
\usepackage{hyperref}
\usepackage{enumitem}
\usepackage{fancyhdr}
\usepackage{lastpage}
\usepackage{xcolor}
\usepackage{verbatim}
\usepackage{caption}
\usepackage{subcaption}
\usepackage{amsmath}
\usepackage{amssymb}
\usepackage{setspace}
\pagestyle{fancy}
\fancyhf{}
\fancyhead[C]{\small\textit{Orchestrate or Upgrade?}}
\fancyfoot[C]{\thepage}
\renewcommand{\headrulewidth}{0.4pt}
\hypersetup{colorlinks=true,linkcolor=blue,urlcolor=blue,citecolor=blue}
\begin{document}
\section{Orchestrate or Upgrade? The Cost-Quality Economics of Multi-Agent Pipelines versus Frontier Models — When Lower Tier Teams Deliver 98\% Quality at 23\% Cost}

\textbf{Youssef Hariri}
DeeWaanAI · osf.io/mjb8x · 2026-08-08

\textbf{Version:} v3 (2026-08-09) — Business of AI edition: business-oriented framing with full academic rigor preserved.

\hrule


\subsection{Abstract}

\hrule


\subsection{Abstract}

Organizations deploying large language models (LLMs) for knowledge work face a
strategic choice: invest in higher-tier models or invest in multi-agent pipeline
structures that orchestrate cheaper models. We test this resource-allocation
tradeoff with a pre-registered 2 (model tier) × 2 (pipeline structure) × 2
(persona) factorial experiment across five LLM families, 38 business-strategy
problems, and 1,900 generation–judgment runs evaluated by a blinded three-judge
panel. Our findings yield a decision framework for AI investment. First, team
structure delivers higher marginal returns at the cheap tier than at the
frontier tier (d = −0.227, p = 0.015): a cheap model with a three-step pipeline
approaches frontier quality at 23\% of the cost (T-EGR = 6.15). Second, personas
degrade output quality on average (d = −0.691), representing a failed investment
in stylistic differentiation. Third, a low-cost refinement pass recovers 60–64\%
of the frontier gap but never achieves paired superiority — a bounded-improvement
result that constrains the substitution thesis. Qualitative coding (κ = 0.824)
reveals that orchestration buys coverage (more structure, more length) while the
frontier wins on density (compressed insight, zero boilerplate). The Pareto
frontier is stable across price assumptions, yielding actionable contingency
rules for budget-constrained firms. We contribute the T-EGR efficiency metric,
a mechanistic account of why structure substitutes for tier only partially, and
a practitioner decision framework that maps budget tier and task type to optimal
pipeline architecture.

\textbf{Keywords:} LLM Orchestration, Multi-Agent Pipelines, Cost-Performance
Analysis, LLM-as-a-Judge, T-EGR, Pre-registration.

\hrule


\subsection{List of Abbreviations}

\begin{table}[htbp]
\centering
\begin{tabular}{ll}
\toprule
Abbr & Definition \\
\midrule
BF₁₀ & Bayes factor in favor of H₁ over H₀ \\
C-S-N & cheap solo, no persona \\
C-S-P & cheap solo, persona \\
C-T-N & cheap team, no persona \\
C-T-P & cheap team, persona \\
F-S-N & frontier solo, no persona \\
F-S-P & frontier solo, persona \\
F-T-N & frontier team, no persona \\
F-T-P & frontier team, persona \\
EGR & Efficiency Gain Ratio \\
FDR & false discovery rate \\
FLOPs & floating-point operations \\
FSN & frontier solo, no persona (baseline) \\
ICER & incremental cost-effectiveness ratio \\
IRR & inter-rater reliability \\
JZS & Jeffreys–Zellner–Siow (Bayes factor prior) \\
LLM & large language model \\
LMM & linear mixed model \\
LOO & leave-one-out \\
OBF & O'Brien–Fleming (spending function) \\
SESOI & smallest effect size of interest \\
T-EGR & Tiered Efficiency Gain Ratio \\
TOST & two one-sided tests (equivalence) \\
\bottomrule
\end{tabular}
\end{table}


\hrule


\section{1 Introduction}

The past two years have seen a rapid industrial shift from single-model
inference to the orchestration of multiple language-model calls: pipelines
chain builder, critic, and refiner agents into structured workflows; agent
teams assign social roles to each step; panels of judges aggregate verdicts
across models. The implicit claim is economic — that *structure can
substitute for model tier*. A practitioner who cannot afford a frontier
model, the argument runs, can reach comparable output quality by composing
cheaper models into an orchestrated system. If true, the claim reshapes how
organizations price, budget, and architect generative-AI products. If
false, orchestration is at best a marginal convenience and at worst a cost
multiplier that never reaches frontier quality.

Prior empirical work on generative AI is encouraging but silent on this
specific claim. Controlled productivity studies show that a single generative
assistant raises output quality and reduces task time (Noy \& Zhang, 2023;
Brynjolfsson, Li, \& Raymond, 2023), yet they treat the model as a given,
single-agent capability. A field experiment with consultants documented a
"jagged technological frontier": large gains inside the model's capability
boundary, large losses outside it (Dell'Acqua et al., 2023). Neither study
varies the structure around the model. Multi-agent and ensemble work shows
that sampling, debate, and tree search can improve accuracy (Wang et al.,
2023; Liang et al., 2024; Yao et al., 2023), but the gains are reported
without cost-matched comparison to a stronger single model. A nascent
theory of test-time compute predicts that, for fixed problem difficulty,
more inference compute on a smaller model can outperform less inference
compute on a larger one (Snell et al., 2024), but its predictions are in
units of FLOPs, not dollars, and only on verifiable tasks.

Three gaps follow. First, no factorial, cost-matched comparison of
orchestrated cheap models versus frontier solo exists on open-ended
tasks. Second, no mechanistic account explains \textit{why} orchestration succeeds
or fails — whether it adds structural coverage or compressed insight.
Third, no economic metric is robust to the price-assumption choices that
inevitably shape per-run costing.

This paper addresses all three. We run a pre-registered 2 (tier) ×
2 (structure) × 2 (persona) factorial experiment over five model families,
38 domain-stratified business problems, and 1,900 generation–judgment
runs, measured by a blinded three-judge panel from families disjoint
from every generator. Cost is stored as tokens and re-priced through a
frozen price vector under three blended input/output ratios, so every
reported dollar is reproducible and re-priceable. A second study applies
a cheap, add-only refinement pass to the orchestrated outputs and codes
the qualitative differences between original, refined, and frontier
answers, testing the mechanistic prediction that orchestration buys
coverage while the frontier wins on density.

Five contributions follow:

\begin{enumerate}
\item \textit{}First cost-matched pre-registered factorial on orchestration
   economics.\textit{} Negative superiority results are reported as primary
   findings, not buried as null results.
\item \textbf{The T-EGR metric} (novelty ratio ÷ cost ratio) and a
   price-assumption-robust Pareto frontier.
\item \textbf{A mechanistic account} from a mixed-methods Study 2 (pooled κ =
   0.824): orchestration buys coverage, the frontier wins on density.
\item \textbf{The refinement boundary}: a \$0.002 refinement pass reaches parity
   with the weakest frontier but never superiority.
\item \textbf{Practitioner contingency rules} mapping cell to use-case to
   expected T-EGR.
\end{enumerate}


The paper proceeds as follows. Section 2 reviews related literature.
Section 3 describes the pre-registered design, apparatus, and analysis,
including an AI-assisted research disclosure (§3.9).
Section 4 presents the results of both studies. Section 5 discusses
theoretical and practical implications. Section 6 acknowledges
limitations. Section 7 outlines future directions. Section 8 concludes.

\hrule


\section{2 Literature Review}

This review synthesizes four streams of research: (i) measured productivity
effects of generative AI, (ii) multi-agent and ensemble generation methods,
(iii) LLM-as-a-judge evaluation methodology, and (iv) open-science methods
for economic measurement. We then integrate these streams with our own prior
work on personas and efficiency metrics to articulate the research gap that
motivates this study.

\subsection{2.1 Productivity effects of generative AI}

\subsubsection{2.1.1 Measured gains in knowledge work}

Three controlled studies establish that generative AI raises knowledge-work
productivity, but they also reveal that the gains are conditional and
bounded. Noy and Zhang (2023) randomized 444 mid-level professionals to
ChatGPT-assisted and control arms on professional writing tasks: assistance
reduced completion time by 0.8 standard deviations and raised output quality
by 0.4 standard deviations, with the largest quality gains among initially
lower-performing workers — a compression of the performance distribution.
Brynjolfsson, Li, and Raymond (2023) observed 5,179 customer-support agents
using a generative assistant in the field and found a 14\% average
productivity increase, concentrated among novice and low-skilled workers
(+34\%), consistent with the model diffusing behaviors of high-ability agents
across the workforce.

A critical limitation unites these studies: both treat the model as a
\textit{given, single-agent capability}. The unit of analysis is one model call,
not an orchestrated system. The professional-writing and customer-support
tasks they study are also relatively closed-form: the desired output is
well-specified (a support response, a draft document), and quality can be
partially verified against ground truth. Open-ended strategic tasks — the
domain of this paper — have no verifiable correct answer, which changes both
the nature of the "frontier" and the applicability of the productivity
gains documented above.

\subsubsection{2.1.2 The jagged frontier and its boundary}

The boundary of these gains is the central caution. Dell'Acqua et al.
(2023) ran a field experiment with 758 consultants on realistic professional
problems and documented a "jagged technological frontier": inside the model's
capability boundary, consultants completed 12.2\% more tasks, 25.1\% faster,
with 40\% higher quality; outside it, consultants using the model were 19
percentage points less likely to produce correct solutions than controls.
Two productive behavioral strategies emerged ("centaurs" dividing work at
the boundary, "cyborgs" interleaving with the model), but neither changed
the location of the boundary itself.

The jagged frontier concept is especially consequential for open-ended tasks.
In closed-form domains (arithmetic, code, support scripts), the boundary is
sharp: a solution is either correct or not. In open-ended business strategy,
the boundary is diffuse — there is no single correct answer, only better or
worse responses. This diffuseness creates a fundamental ambiguity: when a
cheap model produces a plausible but suboptimal strategy, the user may not
know they are "outside the frontier" until the solution fails in deployment.
The open question this paper addresses is whether \textit{orchestration structure}
can shift the effective capability frontier for these open-ended tasks —
and at what cost.

\subsection{2.2 Multi-agent systems, reasoning, and inference-time compute}

\subsubsection{2.2.1 From chain-of-thought to structured reasoning}

The modern reasoning literature begins with the observation that language
models solve complex problems better when they generate intermediate
reasoning steps. Wei et al. (2022) introduced chain-of-thought (CoT)
prompting: appending exemplars that show step-by-step reasoning before the
final answer elicited comparable reasoning in the model, dramatically
improving performance on arithmetic, commonsense, and symbolic reasoning
tasks. The CoT paradigm reframed language-model inference from
single-shot generation to \textit{structured multi-step deliberation} — a
conceptual shift that enabled the orchestration methods reviewed below.

CoT's direct descendant in the accuracy–cost tradeoff space is
self-consistency. Wang et al. (2023) showed that sampling a diversity of
reasoning paths from a single model and taking a majority vote over final
answers improves arithmetic, commonsense, and symbolic reasoning well beyond
CoT decoding alone. The result establishes the general principle that
parallel sampled calls can be converted into accuracy. Its direct corollary
— that accuracy is bought at a multiplicative cost in calls and tokens —
is rarely priced against the alternative of a single stronger model.

\subsubsection{2.2.2 Interactive deliberation: debate, tree search, and agent loops}

Moving beyond independent sampling, recent work has explored \textit{interactive}
structures in which model instances influence each other. Liang et al.
(2024) showed that multiple LLM instances proposing answers, critiquing each
other's responses, and refining their own over debate rounds improves
factuality and reasoning on arithmetic and knowledge benchmarks, encourages
divergent thinking, and reliably converges to shared answers. Yao et al.
(2023) extended the reasoning structure from linear chains to branching
trees: Tree of Thoughts treats each reasoning step as a node in a search
graph, allowing the model to explore multiple reasoning paths, backtrack
from dead ends, and select the most promising branch — deliberate problem
solving rather than single-path generation.

These works demonstrate that structure (sampling, aggregation, critique
loops, tree search) adds quality beyond a single call. What they do not
provide is a cost-matched factorial comparison: prior multi-agent studies
report quality gains without measuring dollars, and without comparing
orchestrated cheap models against an un-orchestrated frontier model on
identical problems. The substitution claim — that structure can stand in
for model tier — is therefore assumed, not tested.

\subsubsection{2.2.3 The test-time compute hypothesis}

A nascent theoretical framework reframes the model-tier-versus-structure
tradeoff as an allocation problem. Snell et al. (2024) demonstrated that
scaling inference-time computation (more tokens, verifier-based search,
iterative refinement) can be more effective than scaling model parameters
for improving performance on reasoning benchmarks. Their compute-optimal
curve predicts that for any fixed problem difficulty, there exists an
optimal allocation between training compute (model size) and inference
compute (test-time search). Below that optimum, a smaller model with more
inference compute outperforms a larger model with less; above it, the
larger model wins.

This framework is the closest theoretical grounding for the "orchestration
economics" claim, but it remains limited in three ways that motivate our
study. First, Snell et al. measure \textit{verifiable} reasoning tasks (math,
code) where a verifier can score intermediate steps; open-ended business
strategy has no verifier. Second, their test-time compute is
\textit{search-based} (sampling + scoring), whereas the multi-agent literature
uses \textit{structural} compute (builder-critic-refiner pipelines, debate rounds)
— the two may have different cost-quality curves. Third, and most
fundamentally, their analysis is in units of FLOPs, not dollars: the
practitioner facing an API bill cannot spend FLOPs. A dollar-matched
empirical test is missing.

\subsubsection{2.2.4 Self-refinement and the limits of iterative improvement}

The final reasoning structure relevant to this study is self-refinement:
a model critiquing and revising its own output. Shinn et al. (2023)
introduced Reflexion, in which a language agent maintains a reflective
memory of past failures and uses verbal reinforcement learning to improve
subsequent trials. On coding and decision-making tasks, Reflexion
outperformed single-shot baselines and even standard CoT by learning from
its own errors.

Self-refinement raises a hypothesis directly tested in our Study 2: if a
cheap model's initial output is below the frontier, can a cheap refinement
pass close the gap? The multi-agent and self-refinement literatures imply
that iterative improvement should raise quality. But they do not tell us
whether the improvement is sufficient to reach parity with a frontier model
that was not itself refined — the marginal comparison that determines
substitution. Nor do they tell us \textit{why} refinement succeeds or fails on
open-ended tasks. Our qualitative coding protocol (Section 3.6.5) is
designed to answer this mechanistic question.

\subsection{2.3 LLM-as-a-judge evaluation methodology}

Because our outcomes are open-ended business solutions, evaluation itself is
a design object. Zheng et al. (2023) established LLM-as-a-judge as a
scalable evaluation paradigm, showing strong agreement between strong judge
models and human preferences on pairwise comparison tasks, while documenting
three systematic biases that threaten validity: position bias (preferring
the candidate presented first), verbosity bias (preferring longer outputs),
and self-enhancement bias (preferring outputs from the same model family as
the judge). Their mitigations — swapping candidate order, controlled
few-shot examples, length normalization — inform our blinding and
order-shuffling protocol.

The self-enhancement bias is particularly acute for multi-agent studies. If
the generator and judge share model architecture, the judge may reward the
generator's stylistic conventions rather than substantive quality. Verga et
al. (2024) extended the single-judge design into a jury: a Panel of LLM
evaluators (PoLL) composed of several smaller models from different families
correlated better with human judgments than a single large judge, at lower
cost, while mitigating the self-preference bias that afflicts same-family
evaluation. Our evaluation framework inherits this logic directly: three
judges from families disjoint from all generators, each under three frozen
prompt variants, scoring blind to cell and family. This design is especially
important for detecting the "persona penalty" hypothesized in this study —
a stylistic shift that a same-family judge might reward while a
disjoint-family jury penalizes.

\subsection{2.4 Open-science methods and economic measurement}

The fourth stream supplies the discipline this paper applies. Nosek et al.
(2018) argue that preregistration distinguishes confirmatory from
exploratory analysis, converting transparency into credibility; our ten
protocols and their amendments were all registered before the corresponding
computations ran. Steegen et al. (2016) show that conclusions can vary
across reasonable analytic choices and recommend reporting a multiverse of
analyses; our price-assumption panel (three blended input/output ratios)
and pre-registered robustness battery implement this idea. Sandve et al.
(2013) codify reproducible computational research — versioned scripts,
recorded steps, stable environments — which our token ledger and frozen
configuration files follow.

Two statistical practices complete the foundation. Albers and Lakens
(2018) demonstrate that effect sizes estimated from pilots are unstable and
bias follow-up power analyses, so pilot data must never anchor confirmatory
claims; accordingly, our pilot informed only feasibility gates (G1/G2), never
the main analysis. Lakens (2013) makes the case for reporting effect sizes
with confidence intervals to support cumulative science; we report Welch's d
with bootstrap intervals throughout. Cohen's (1960) κ coefficient provides
our inter-rater and coder-agreement metric.

Finally, measuring the economics of inference requires separating quantity
from price. Diewert (1976) formalized the decomposition of value changes
into price and quantity components in index-number theory; we apply the same
logic by storing tokens (quantity) and multiplying by a frozen price vector,
so every reported dollar is reproducible and re-priceable. Tomlinson et al.
(2024) quantify the environmental footprint of generative tasks, reinforcing
that inference cost — economic and environmental — is a first-order design
variable rather than an implementation detail.

\subsection{2.5 Integration: from persona economics to orchestration economics}

Our own prior work supplies the immediate empirical background. In a
cross-architecture evaluation of cultural personas across seven large
language models, persona effectiveness followed a contingency model rather
than a universal rule, and the Efficiency Gain Ratio (EGR) showed that small
models with personas achieved superior performance at 40–88\% lower cost
(Hariri, 2026b). The same line of work measured creativity with multi-judge
panels over thousands of generated business solutions, establishing
Novelty as the best-agreed creativity dimension (Hariri, 2025a), and showed
that team composition effects are contingent on task verifiability
(Hariri, 2025b).

Integrating the four streams above with this lineage yields a precise gap
and three pre-registered hypotheses, each grounded in the literature:

\textbf{Substitution hypothesis (H1):} The test-time compute framework (Snell et
al., 2024) predicts that structure can substitute for model tier at the
cheap end of the compute-optimal curve. We test this prediction with a
factorial tier × structure design, priced in dollars, on open-ended tasks
where no verifier exists. If the hypothesis holds, cheap teams should
approach frontier-solo quality at a fraction of the cost.

\textbf{Persona hypothesis (H2):} Our prior work (Hariri, 2026b) showed that
cultural personas shift quality within single-model cells. We extend this
to teams: if personas amplify structural rigor, the persona effect should be
larger for team structures than for solo structures. If personas primarily
add stylistic variation without substance, the persona-loaded cells should
be Pareto-dominated by their no-persona twins.

\textbf{Refinement hypothesis (H3):} The self-refinement literature (Shinn et
al., 2023) suggests that a cheap critique-and-revise pass can recover
quality gaps. We test whether a frozen, add-only refiner can close the
residual gap between orchestrated cheap outputs and frontier outputs — and
whether the mechanism is additive coverage (more structure) or compressed
insight (better density).

The experiment in Section 3 tests all three hypotheses, priced. The
qualitative coding protocol (Section 3.6.5) is designed to distinguish the
two mechanisms implied by H3: "orchestration buys coverage, the frontier
wins on density" is the falsifiable prediction derived from the literature
review above.

\hrule


\section{3 Methodology}

\subsubsection{3.1 Research Design}

The experiment is a 2 × 2 × 2 between-cell factorial with model
\textit{family} as a blocking factor. The three manipulated factors are
\textbf{tier} (frontier flagship vs. cheap sibling within the same
generation), \textbf{structure} (solo single call vs. three-step team of
builder → critic → refiner), and \textbf{persona} (none vs. a fixed
cultural-persona mapping applied to each team role). Five model families
— Kimi, Qwen, GPT-5.6, Claude, and Gemini — supply both frontier and
cheap tiers. Thirty-eight business problems, stratified across
knowledge-work domains, are identical across all cells; no instruction
wrappers are added to the rendered prompt. Every design and statistical
decision is time-stamped, numbered (D1–D39), and registered in a
frozen decision log; ten protocols and their amendments were registered
before the corresponding computation ran, maintaining the confirmatory/
exploratory separation required by preregistration discipline (Nosek et
al., 2018).

\begin{figure}[htbp]
\centering
\includegraphics[width=\textwidth]{/Users/youssefhariri/Documents/app 25/research/orchestration-economics/paper/figures/fig1_design_schematic.png}
\caption{Figure 1. Experimental Design: 2 × 2 × 2 Factorial + Apparatus}
\end{figure}


\subsubsection{3.2 Models Evaluated}

\paragraph{3.2.1 Frontier Models (Five Families)}

Each family contributes its current \textit{mainline} flagship, pinned on
2026-07-30: \texttt{kimi-k3}, \texttt{qwen3.7-max-2026-06-08}, \texttt{gpt-5.6-sol},
\texttt{claude-opus-5}, and \texttt{gemini-3.1-pro-preview}. Above-mainline tiers
(e.g. \texttt{claude-fable-5}) are deliberately excluded because they would
break cross-family comparability. Gemini's preview status carries a
pre-registered fallback to \texttt{gemini-2.5-pro} if the model is retired
mid-experiment; the fallback was never invoked. All models run with
generator temperature = 1, a uniform value forced by the OpenAI and
Moonshot APIs (which reject other values) and applied identically
across families for ceteris-paribus comparability (D27).

\paragraph{3.2.2 Cheap Siblings}

Each family contributes one cheap-tier model from the same generation:
\texttt{kimi-k2.6}, \texttt{qwen3.7-plus-2026-05-26}, \texttt{gpt-5.6-luna},
\texttt{claude-haiku-4-5-20251001}, and \texttt{gemini-3.1-flash-lite}. Price ratios
(frontier : cheap) range from 3.4× to 8× across families. Family-level
heterogeneity in price ratio is absorbed by the family random effect
in our mixed-model analysis (§3.6.2); it is not a confound of the
primary tier × structure test.

\subsubsection{3.3 Orchestration Structures and Personas}

\textbf{Solo} runs invoke a single model call on the rendered problem.
\textbf{Team} runs execute a three-step pipeline — builder produces a draft,
critic critiques it, refiner produces the final answer — with all
intermediate outputs discarded from the judged answer (and therefore
from the quality score) but retained in the cost ledger (a disclosed
cost bias, favorable to the cheap-team cells). Persona is implemented
as a fixed mapping from role (builder/critic/refiner) to a cultural
persona description, identical across all problems and families
(Hariri, 2026b). The mapping is pre-registered and never varied
per-problem; it is therefore a treatment level, not a confound.

\subsection{3.4 Stimuli and response generation}

Thirty-eight domain-stratified business problems are used. Fifteen are
inherited verbatim from a prior validated set (Hariri, 2025a; DOI
10.5281/zenodo.17407392); ten were drafted post-cutoff to control for
contamination and validated as structurally identical; the remainder
extend the grid to cover five domains equally. Each problem is rendered
as \texttt{background\_context + "\n\n" + business\_problem} with no
instructional wrapper — adding wrapper text would introduce an
unmeasured prompt treatment. Generator temperature is fixed at 1
(§3.2.1). Every run is time-stamped and flagged for whether it executed
inside the pre-registered preferred window (14:00–00:00 UTC), which
simultaneously captures the Qwen off-peak discount and fully avoids
both DeepSeek peak-surge windows; window-free providers (Kimi,
Together-hosted models, Western providers) are time-invariant and
therefore unaffected.

\subsubsection{3.5 Evaluation Framework}

\paragraph{3.5.1 Judge Model Selection}

Three judges from families disjoint from every generator:
\texttt{deepseek-chat} (resolved to \texttt{deepseek-v4-flash}), \texttt{zai-org/GLM-5.2},
and \texttt{meta-llama/Llama-3.3-70B-Instruct-Turbo}. The disjointness rule
(D33) guards against self-enhancement bias (Zheng et al., 2023): even
when deepseek is the registered fallback generator, the fallback was
never invoked, so disjointness holds in the executed data.

\paragraph{3.5.2 Scoring Rubric}

Six 1–5 dimensions are scored: \textit{}Novelty, Usefulness, Flexibility,
Elaboration, Cultural Sensitivity, and an aggregate Practicality\textit{}.
The composite primary outcome is the mean of the first five
dimensions. Three prompt variants are used per judge, frozen verbatim
and never re-optimized after a prior judge-prompt optimization attempt
collapsed κ in earlier work (Hariri, 2025a). Every solution is
scored by all three judges under all three variants (nine scores per
solution), with candidate order shuffled per variant to neutralize
position bias (Zheng et al., 2023). Judges run at temperature 0 for
determinism.

\paragraph{3.5.3 Inter-Rater Reliability}

Pairwise judge correlations on per-run composite scores are 0.87
(deepseek–zhipu), 0.61 (deepseek–meta), and 0.49 (meta–zhipu). The
panel-level reliability exceeds any single judge's correlation with
the latent quality; this mirrors the jury logic of Verga et al.
(2024) in which diversity of smaller judges outperforms a single
large judge.

\subsubsection{3.6 Statistical Analysis}

\paragraph{3.6.1 Effect Sizes}

Welch's d with bootstrap 95\% confidence intervals (5,000 resamples,
seed 42) throughout (Lakens, 2013). Estimation-first reporting:
intervals, not p-values, carry the weight of the argument; p-values
are reported for completeness and to comply with sequential-analysis
boundaries (§3.6.2).

\paragraph{3.6.2 Inferential Statistics}

The primary test of H1 is the pre-registered interaction contrast
\texttt{d = (FT−FS) − (CT−CS)}, expected negative (team helps cheap more than
frontier). Inference uses a mixed model with random effects for
\textit{problem} and \textit{family}, solving the pseudo-replication identified in
the author's prior work (Hariri, 2026b). Sequential analysis with an
O'Brien–Fleming spending function (protocol 03) governs interim looks
and the final stopping boundary; the Bayesian JZS Bayes factor is
reported as recommended (A2.4) but does not gate the decision.
Family-wise corrections within secondary hypotheses use
Benjamini–Hochberg FDR (α = 0.05); Holm is used for the small set of
directional family comparisons in Study 2.

\paragraph{3.6.3 Economic Analysis}

Cost is stored as tokens in a run-level ledger and re-priced through a
frozen price vector (\texttt{config/price\_vector.json}, version 2026-07-30).
Because the ledger does not split input from output tokens, cost is
computed under a blended price
\texttt{blend(φ) = φ · p\_out + (1 − φ) · p\_in} with a primary φ = 0.50 and a
sensitivity panel φ ∈ \{0.25, 0.50, 0.75\} (Diewert, 1976, provides the
price–quantity decomposition precedent; Steegen et al., 2016, the
multiverse-robustness precedent). Frontier-membership claims must be
stable across all three φ. The \textit{}Tiered Efficiency Gain Ratio
(T-EGR)\textit{} is \texttt{(Quality\_cell / Quality\_F-S-N) / (Cost\_cell / Cost\_F-S-N)},
computed within each family and pooled; F-S-N is the baseline.
\textbf{ICER} is reported relative to the cheapest cell. Judge cost is
excluded from unit T-EGR (it is shared infrastructure, not cell
compute; a disclosed exclusion).

\paragraph{3.6.4 Robustness Checks}

Pre-registered battery (A2): leave-one-family-out, leave-one-
problem-out, mean-vs-median, BH-FDR on secondary dimensions,
length-confound adjustment on H1, and empirical-Bayes partial
pooling of per-family interaction estimates. Random-effects
meta-analysis across the five families confirms the interaction
generalizes (pooled d = −0.270, I² = 28.5\%).

\paragraph{3.6.5 Qualitative Coding Protocol (Study 2)}

Thirty triples (original CTP / refined CTP / same-problem same-family
FSN) are sampled deterministically (seed 20260801, 10 per family).
Two blinded LLM coders perform a two-wave coding: free induction on
a subset, then closed coding against a frozen 13-code codebook
(six RQ1 insight codes, six RQ2 structure codes, one RQ3 bloat code).
Cohen's κ gate is 0.70; codes below threshold are downgraded to
medium reliability and disclosed (Cohen, 1960). The prompt presented
to coders contains the three texts only — no cell label, family, or
score, preserving blind evaluation.

\subsubsection{3.7 Validity and Reliability}

\textbf{Construct validity:} the rubric separates Novelty from Usefulness/
Elaboration, so a stylistic shift (e.g. more discourse markers from a
persona) is not conflated with substantive quality. \textit{}Internal
validity:\textit{} frozen configurations, blinded judging, shuffled order,
disjoint judge families. \textbf{External validity:} single task domain
(English business strategy) is the acknowledged boundary. \textit{}Disclosed
threats:\textit{} intermediate team outputs excluded from judged answers bias
team costs downward (favorable to cheap teams, disclosed not
corrected); truncation at source affects a small set of long runs and
is audited separately.

\subsubsection{3.8 Code and Data Availability}

All analysis code, the frozen token ledger, judge rationales, and the
ten protocols are committed to the \texttt{orchestration-economics}
repository; the OSF registration is \texttt{osf.io/mjb8x}. A reproduction
runbook (Sandve et al., 2013) lists the commands that regenerate every
number reported in §4.

\subsection{3.9 AI-assisted research disclosure}

This study used large language models as research tools under continuous
human supervision. Specifically, \textbf{Kimi} (Moonshot AI, \texttt{kimi-k3}) and
\textbf{Qwen} (Alibaba Cloud, \texttt{qwen3.7-max}) were employed in the following
capacities: (1) \textbf{ideation} — brainstorming experimental design
alternatives, hypothesis formulations, and analysis strategies during
the protocol-drafting phase; (2) \textbf{code generation} — producing initial
versions of the execution pipeline, statistical analysis scripts, and
figure-generation code, all of which were reviewed, corrected, and
approved by the author before execution; (3) \textbf{audit} — reviewing
citation accuracy, checking statistical assumptions, and identifying
inconsistencies across protocol documents; (4) \textbf{execution} — running
the generation and judging campaigns through API calls, with all
configurations frozen and logged before each run.

No LLM made unsupervised decisions about the experimental design,
statistical analysis, or the interpretation of results. All ten
pre-registered protocols, all 39 design decisions (D1–D39), and all
analysis choices were authored, reviewed, and approved by the human
author. LLM-generated code was treated as a first draft, not a final
product: every script was inspected, every output was validated against
expected ranges, and every reported number traces to a frozen,
version-controlled analysis pipeline. The qualitative coding protocol
(§3.6.5) used two LLM coders with blinded evaluation — a design choice
that converts the LLM's systematic biases into a measured inter-rater
reliability coefficient (κ = 0.824) rather than an unmeasured confound.

The use of LLMs as evaluation judges (§3.5) is a distinct methodological
choice from the use of LLMs as research tools: the judges evaluate
\textit{generated business solutions}, not the research process itself. The
disjointness rule (D33) ensures that no judge shares architecture with
any generator, mitigating the self-enhancement bias documented by
Zheng et al. (2023).

\hrule


\section{4 Results}

\subsubsection{4.1 The Interaction Effect: Does Pipeline Structure Pay Off More for Cheap Models?}

From a resource-allocation perspective, the most important finding is that
pipeline structure delivers higher marginal returns at the cheap tier than at
the frontier tier. The pre-registered interaction contrast is d = −0.227
(bootstrap 95\% CI [−0.410, −0.040], Welch t = −2.45, p = 0.0145), with the
sequential-analysis boundary crossed at the scheduled final look. For a firm
making an AI budget decision, this implies that pipeline dollars are better
spent upgrading cheap models to team configurations than adding team overhead
to already-expensive frontier models. The JZS Bayes factor for the interaction
is 1.002 (inconclusive), reported as recommended but not gating the decision.

The secondary effects provide additional strategic context. The \textit{}tier
main effect\textit{} (d = +0.294, CI [+0.201, +0.382], BF₁₀ = 1.7 × 10⁷) confirms
that frontier models command a quality premium, but the magnitude is dwarfed
by the cost difference. The \textbf{structure main effect} (d = +0.117, CI
[+0.030, +0.205]) shows that teams slightly outperform solos on average — a
small but consistent productivity gain from orchestration. The \textbf{persona main effect} is large and in the \textit{opposite}
direction to the pre-registered expectation: d = −0.691 (CI [−0.761,
−0.620], BF₁₀ = 1.0 × 10⁵⁰). Persona-loaded cells score \textit{lower} than
their no-persona twins, representing a failed investment in stylistic
differentiation. From a business perspective, personas increase cost while
degrading output quality — a value-destroying configuration. A post-hoc dimension decomposition shows that
Novelty is the sole dimension where personas help (+0.154, p\_FDR =
0.001) while Usefulness (−1.024) and Elaboration (−0.966) drive the
reversal. The structure × persona interaction is not significant (F =
1.23, p = 0.27).

\textbf{Strategic Implication:} Pipeline structure is a cheap-tier amplifier,
not a universal substitute. For budget-constrained firms, investing in team
pipelines for cheap models yields the highest marginal return. Personas, by
contrast, represent value destruction — they increase cost while degrading
quality.

\begin{figure}[htbp]
\centering
\includegraphics[width=\textwidth]{/Users/youssefhariri/Documents/app 25/research/orchestration-economics/paper/figures/fig21_sequential_boundary.png}
\caption{Figure 21. Sequential analysis with O'Brien–Fleming spending function. The test statistic z = −2.447 crosses the interim-2 boundary |z| ≥ 1.977 (dashed line). The shaded region is the continuation region. The pre-registered decision is to stop and report H1 as supported.}
\end{figure}


\begin{figure}[htbp]
\centering
\includegraphics[width=\textwidth]{/Users/youssefhariri/Documents/app 25/research/orchestration-economics/paper/figures/fig2_forest_plot.png}
\caption{Figure 2. Quality by Experimental Cell with 95\% Confidence Intervals}
\end{figure}


\subsubsection{4.2 The Hidden Cost of Personas: Style Without Substance}

Every persona-loaded cell is Pareto-dominated by its no-persona twin
on the cost–quality plane: at equal or lower cost, the N cell scores
higher on composite quality. The penalty is largest on dimensions that
rubric readers associate with substance (Usefulness, Elaboration),
smallest on the one most susceptible to stylistic signaling (Novelty).
This pattern mirrors the contingency finding in the author's prior
work (Hariri, 2026b) and is consistent with the jagged-frontier
mechanism (Dell'Acqua et al., 2023): personas push the generator
toward a more stylized region of its output space that the judge,
correctly, penalizes on substance.

\textbf{Strategic Implication:} Personas degrade output quality while increasing
cost — they are value-destroying for knowledge-work tasks. Firms should treat
them as a failed investment rather than a quality lever.

\begin{figure}[htbp]
\centering
\includegraphics[width=\textwidth]{/Users/youssefhariri/Documents/app 25/research/orchestration-economics/paper/figures/fig22_persona_radar.png}
\caption{Figure 22. Quality-dimension radar for persona-loaded (red, dashed) versus no-persona (blue, solid) cells. The persona polygon is visibly smaller on Usefulness and Elaboration, slightly larger on Novelty — the visual signature of the style-over-substance tradeoff.}
\end{figure}


\begin{figure}[htbp]
\centering
\includegraphics[width=\textwidth]{/Users/youssefhariri/Documents/app 25/research/orchestration-economics/paper/figures/fig10_persona_dimensions.png}
\caption{Figure 10. Persona Effect by Quality Dimension (FDR-Corrected)}
\end{figure}


\subsubsection{4.3 The Frontier Gap: Can Cheap Teams Match Frontier Performance?}

Paired comparisons between cheap-team-with-persona (CTP) and the
frontier-solo-no-persona baseline (F-S-N) are negative in all five
families (kimi d = −0.71, claude d = −1.42, qwen d = −2.17; all CIs
exclude zero after Holm correction). The gap is the empirical target
that a refinement pass must close — the question of Study 2.

\textbf{Strategic Implication:} As executed, cheap teams do not beat a single
frontier call on matched problems. The substitution thesis is bounded:
structure improves cheap models but cannot fully close the frontier gap
without additional investment.

\subsubsection{4.4 The Refinement Play: Closing the Gap at Low Marginal Cost}

A frozen, add-only refiner is applied to the CTP outputs. The pilot
satisfied the feasibility gate (G1, 100\% completion); the full-pool
extension covered 171 pairs per family. The pre-registered superiority
gate (G2: d ≥ 0.2 with CI excluding 0 in all three families, in the
direction refined > FSN) is mechanically failed in all three
families: kimi d = −0.295, claude d = −1.301, qwen d = −0.809 (all
CIs strictly negative after Holm). The superiority claim is formally
rejected.

Refinement nevertheless does substantial repair work. Measured as gap
recovery against each family's own pre-refinement CTP-vs-FSN gap, the
pass recovers 60\% of the gap for kimi, 8\% for claude (whose pre-
refinement gap was already small), and 64\% for qwen. The Novelty
protection criterion (H2.2: no family loses more than −0.10 Novelty)
is passed by all three; Novelty actually \textit{increases} for qwen
(d = 1.10, p\_FDR < 0.001). Usefulness and Elaboration show large
gains in kimi and qwen (d = 0.8 to 1.4, all p\_FDR < 0.001) and are
unchanged in claude — which was already at ceiling on those
dimensions before refinement.

Length analysis confirms the dilution hypothesis is \textit{not} the driver:
the pooled Pearson correlation between length ratio (refined/original)
and composite-score change is r = +0.68, Spearman ρ = +0.71
(p < 0.001). Longer refined answers score higher, not lower.

Cross-family comparisons (Figure 7) remain negative after refinement:
no cross-family CTP pairing flips to superiority, confirming that the
within-family result is not an artifact of family selection.

\textbf{Strategic Implication:} A \$0.002 refinement pass delivers bounded
improvement — it recovers most of the frontier gap on average but never
achieves paired superiority. For firms, this means refinement is a viable
risk-mitigation strategy (improve quality cheaply) but not a replacement
strategy (it will not beat the frontier).

\begin{figure}[htbp]
\centering
\includegraphics[width=\textwidth]{/Users/youssefhariri/Documents/app 25/research/orchestration-economics/paper/figures/fig4_gap_recovery.png}
\caption{Figure 4. Study 2: Refinement Probe Results}
\end{figure}


\begin{figure}[htbp]
\centering
\includegraphics[width=\textwidth]{/Users/youssefhariri/Documents/app 25/research/orchestration-economics/paper/figures/fig7_cross_family_matrix.png}
\caption{Figure 7. Cross-Family Comparison Matrix (Study 2)}
\end{figure}


\subsubsection{4.5 Why Orchestration Buys Coverage but Not Density}

Thirty triples (10 per family) are coded under a 13-code codebook.
Pooled Cohen's κ = 0.824 (93.1\% agreement), well above the 0.70
gate. Six of the thirteen codes fall below 0.70 (diagnostic\_reframe
0.54, named\_concepts 0.49, assumption\_test 0.56, structured\_report
0.59, external\_authority 0.56, redundancy\_bloat 0.65), a base-rate
artifact: percentage agreement on these codes ranges from 82\% to 98\%,
but κ is depressed by near-zero prevalence. Conclusions on
below-threshold codes are disclosed as medium reliability.

Of the six codes whose prevalence changed between original and
refined (kimi, qwen), all six are \textit{structural scaffolding} codes:
operational\_parameters, structured\_report, phased\_timeline,
tables\_matrices, quantified\_targets, granular\_ops\_detail. No insight
code (diagnostic\_reframe, named\_concepts, assumption\_test,
concise\_maxims) moved. The convergence metric shows refinement does
not consistently move cheap outputs toward the frontier on insight
codes (kimi toward 4 / away 6; qwen 5/4; claude 0/2 — claude was
already at ceiling).

The residual frontier edge is precisely located. The only
κ-reliable frontier-vs-refined difference ≥ 0.30 is
\texttt{concise\_maxims}: qwen frontier 1.00 vs refined 0.55 (gap 0.45).
The medium-reliability \texttt{redundancy\_bloat} code is 0.00 in every
frontier sample and 0.65–0.95 in every refined sample. Claude's
original answers are already at the codebook ceiling (≥ 0.90 on 11
of 13 codes); the number of changed codes between original and
refined is zero — a saturation mechanism that explains the 8\% gap
recovery in §4.4.

\textbf{Strategic Implication:} Orchestration is a coverage play — it adds
structure, length, and completeness. The frontier's advantage is a density
play — compressed insight with zero waste. Add-only refinement cannot
manufacture density, which explains why cheap teams approach but do not
surpass frontier quality.

\begin{figure}[htbp]
\centering
\includegraphics[width=\textwidth]{/Users/youssefhariri/Documents/app 25/research/orchestration-economics/paper/figures/fig5_code_heatmap.png}
\caption{Figure 5. Code Prevalence Heatmap: Original, Refined, and Frontier (Study 2)}
\end{figure}


\subsubsection{4.6 The Pareto Frontier: Optimal Configurations for Budget-Constrained Firms}

The cost–quality map is shown in Figure 3. At the primary φ = 0.50,
the cell means are:

\textbf{Table 1.} Cell summary statistics (primary analysis, φ = 0.50).

\begin{table}[htbp]
\centering
\begin{tabular}{lllll}
\toprule
Cell & n & Cost (USD) & Quality & T-EGR \\
\midrule
C-S-P & 285 & 0.0145 & 4.210 & 8.60 \\
C-S-N & 285 & 0.0194 & 4.507 & 6.15 \\
C-T-P & 285 & 0.0715 & 4.327 & 1.58 \\
C-T-N & 285 & 0.0734 & 4.531 & 1.49 \\
F-S-P & 190 & 0.0670 & 4.406 & 1.36 \\
F-S-N & 190 & 0.0835 & 4.586 & 1.00 (base) \\
F-T-N & 190 & 0.2951 & 4.610 & 0.27 \\
F-T-P & 190 & 0.3044 & 4.370 & 0.27 \\
\bottomrule
\end{tabular}
\end{table}


\begin{figure}[htbp]
\centering
\includegraphics[width=\textwidth]{/Users/youssefhariri/Documents/app 25/research/orchestration-economics/paper/figures/fig3_cost_scatter.png}
\caption{Figure 3. Cost-Effectiveness Frontier: 8 Cells + Study-2 Refined CTP}
\end{figure}


\begin{figure}[htbp]
\centering
\includegraphics[width=\textwidth]{/Users/youssefhariri/Documents/app 25/research/orchestration-economics/paper/figures/fig6_tegr_bars.png}
\caption{Figure 6. Tiered Efficiency Gain Ratio (T-EGR) by Cell}
\end{figure}


The Pareto frontier is \textbf{identical across all three φ}: \{C-S-P,
C-S-N, C-T-N, F-S-N, F-T-N\}. The dominated set \{C-T-P, F-S-P, F-T-P\}
contains every persona-loaded team cell — the economic echo of the
§4.2 persona penalty. Frontier membership is stable across the
sensitivity panel, satisfying the pre-registered claim rule.

\begin{figure}[htbp]
\centering
\includegraphics[width=\textwidth]{/Users/youssefhariri/Documents/app 25/research/orchestration-economics/paper/figures/fig19_phi_sensitivity.png}
\caption{Figure 19. Pareto frontier stability across the price-assumption sensitivity panel. Left: φ = 0.25 (input-heavy blend). Center: φ = 0.50 (primary analysis). Right: φ = 0.75 (output-heavy blend). The frontier set \{C-S-P, C-S-N, C-T-N, F-S-N, F-T-N\} is identical in all three panels, satisfying the pre-registered cross-φ stability rule.}
\end{figure}


ICER (vs. the cheapest cell, C-S-P): C-S-N is the best buy at
+\$0.016 per point of composite quality; F-T-P is the worst at
+\$1.81 per point. The refined kimi C-T-P point (not in the table)
sits at \$0.0654 / quality 4.585 — absolute parity with pooled F-S-N
(4.586) at 78\% of its cost, with the caveat that the paired
within-family comparison remains negative (disclosure per §4.4).

\textbf{Strategic Implication:} The economically optimal strategy is to
substitute structure for tier at the cheap end (C-S-N: 98\% of frontier
quality at 23\% of the cost). The worst dollar in the experiment is spent
on frontier teams — they cost 3.5× the frontier-solo price for negligible
quality gain.

\subsubsection{4.7 Robustness: Do the Results Hold Across Model Families?}

The A2 battery supports the §4.1 decision. Leave-one-family-out: d
ranges from −0.143 (drop kimi) to −0.303 (drop gemini), sign stable
in all five folds; no single family drives the interaction. Leave-
one-problem-out: d ranges from −0.212 to −0.247 across all 19
folds; no problem drives the interaction. Sensitivity power at the
smallest cell (n = 190) detects d = 0.288 at 80\% power; the
observed |d| = 0.227 is detected through the within-subjects
contrast over all 1,900 runs rather than per-cell power (A2.7).
Length adjustment on the H1 contrast yields p = 0.63 — the
interaction is not a length confound. The random-effects meta-
analysis over the five families (A3.4) gives a pooled d = −0.270
(CI [−0.472, −0.069], z = −2.64, p = 0.008) with low-moderate
heterogeneity (τ² = 0.015, I² = 28.5\%), confirming the
interaction generalizes.

Results are stable (mean absolute change of cell composite under
the battery < 0.04 across perturbations). The leave-one-out
sensitivity (Figure 8) and random-effects meta-analysis (Figure 9)
confirm that no single family, problem, or outlier drives the result.

\textbf{Strategic Implication:} The findings are robust to judge, model
family, problem domain, and output-length variations. The interaction is
small but reliable and generalizes across the five leading LLM families —
a stable result for firms making platform-agnostic investment decisions.

\begin{figure}[htbp]
\centering
\includegraphics[width=\textwidth]{/Users/youssefhariri/Documents/app 25/research/orchestration-economics/paper/figures/fig8_robustness_loo.png}
\caption{Figure 8. Robustness: Leave-One-Out Sensitivity (A2 Battery)}
\end{figure}


\begin{figure}[htbp]
\centering
\includegraphics[width=\textwidth]{/Users/youssefhariri/Documents/app 25/research/orchestration-economics/paper/figures/fig9_meta_analysis_forest.png}
\caption{Figure 9. Random-Effects Meta-Analysis: Interaction by Family}
\end{figure}


\subsubsection{4.8 Summary of Key Findings}

\textbf{Table 2.} Summary of key findings.

\begin{table}[htbp]
\centering
\begin{tabular}{lll}
\toprule
Finding & Evidence & Section \\
\midrule
Team structure helps cheap more than frontier & H1 contrast d = −0.227 & §4.1 \\
Frontier > cheap overall & tier main effect d = +0.294 & §4.1 \\
Persona penalty on net & persona main effect d = −0.691 & §4.1 \\
Persona-loaded cells Pareto-dominated & cost–quality plane & §4.2 \\
CTP < FSN in every family & paired comparisons & §4.3 \\
Refinement recovers 60–64\% of gap (kimi, qwen); G2 FAIL & H2.1/H2.3 & §4.4 \\
Orchestration buys coverage; frontier wins on density & κ = 0.824; 13 codes & §4.5 \\
Pareto frontier stable across φ & \{C-S-P, C-S-N, C-T-N, F-S-N, F-T-N\} & §4.6 \\
\bottomrule
\end{tabular}
\end{table}


\hrule


\section{5 Discussion}

\subsubsection{5.1 A Decision Framework for AI Investment}

The results support a contingency framework for AI investment rather
than a universal rule. For budget-constrained firms, \textit{}invest in cheap-team
pipelines for coverage-sensitive tasks\textit{} (timelines, matrices, quantified
targets): C-S-N achieves 98\% of frontier quality at 23\% of the cost
(T-EGR 6.15) — the highest-return configuration in the study. \textit{}Do not
orchestrate at the frontier tier\textit{}: frontier team cells spend 3.5× the
frontier-solo cost for +0.024 points of quality (T-EGR 0.27), making them
the worst dollar in the experiment. \textit{}Use a cheap refinement pass as a
risk-mitigation strategy\textit{} when absolute quality parity is the target:
the refined kimi point reaches parity at 78\% of F-S-N's cost, the cheapest
route to that milestone. \textbf{Avoid personas}: they increase cost while
degrading quality — a value-destroying configuration.

\begin{figure}[htbp]
\centering
\includegraphics[width=\textwidth]{/Users/youssefhariri/Documents/app 25/research/orchestration-economics/paper/figures/fig20_contingency_rules.png}
\caption{Figure 20. Practitioner contingency rules. Decision flowchart mapping budget tier and task type to the recommended orchestration cell, with expected T-EGR. Derived from the §5.1 contingency model.}
\end{figure}


\subsubsection{5.2 The Density Ceiling: Why Additive Pipelines Hit a Wall}

The qualitative mechanism (§4.5) explains the pattern. Orchestration
is compositional — add a critic, add a refiner, add a section, add a
table. Composition grows the answer but cannot compress it. The
frontier's residual advantage sits exactly where additive composition
fails: \texttt{concise\_maxims} (compressed, quotable strategic principles)
and zero \texttt{redundancy\_bloat}. These are the outputs of a single model
that has been selected or trained to produce dense prose. An add-only
refiner, by design, cannot delete the boilerplate the team structure
produced. The falsifiable prediction that follows is that a
\textit{subtractive} refiner — one permitted to cut as well as add — would
narrow the density gap. This is the central future-work experiment
implied by our results.

\subsubsection{5.3 Contribution to the Business of AI Literature}

Our findings extend the author's prior lineage in two directions.
The EGR metric introduced in Hariri (2026b) measured the cost
performance of \textit{personas in single models}; T-EGR extends it to
teams and pipelines, preserving the ratio form while re-basing on
the frontier cell. The team-composition contingency identified on
verifiable tasks in Hariri (2025b) reappears here on open-ended
tasks, but with the sign reversed for personas (harmful rather than
neutral on substance). The jagged-frontier finding of Dell'Acqua et
al. (2023) is replicated as a \textit{cell-level} result: cheap models
inside the persona-free, coverage-oriented region beat their cost-
matched frontier alternatives; cheap models pushed outside by
persona-loading or by density demands fall below the frontier.

The test-time-compute prediction of Snell et al. (2024) is partially
confirmed: more inference compute (team, refinement) improves the
cheaper model, but only up to a parity ceiling, and only when priced
in dollars does the tradeoff reveal its practitioner shape. Our
contribution is the dollarization of the compute-optimal curve on
open-ended tasks where no verifier exists.

\subsubsection{5.4 Methodological Lessons for AI Deployment Research}

Two lessons carry beyond this study. First, preregistered negative
superiority results can be the spine of a paper, not its appendix.
The G2 mechanical failure in §4.4 is as informative as any positive
finding, and preregistration is what made it reportable as a primary
result rather than as a null. Second, token-invariant costing —
storing tokens and re-pricing through a frozen vector — is reusable
infrastructure for the field: every reported dollar is reproducible
by any reader who supplies an alternative price vector (Diewert,
1976; Sandve et al., 2013).

\hrule


\section{6 Limitations}

Four acknowledged limitations bound the claims. (1) Evaluation is
LLM-only; no human panel was recruited. Our bias controls (blinding,
order shuffle, disjoint judge families, three prompt variants) follow
the best current practice (Zheng et al., 2023; Verga et al., 2024)
but do not substitute for human agreement on a subset. (2) The
blended-φ price assumption is a disclosed limitation; we mitigate it
with a pre-registered sensitivity panel and require cross-φ stability
for any frontier-membership claim. (3) The task domain is English
business strategy; generalization to code, scientific writing, or
non-English domains is an empirical question. (4) Cross-family
comparisons (e.g. kimi vs. claude refinement rates) are descriptive
and confounded with family quality; only within-family comparisons
support causal language. Two further disclosures — team intermediate
outputs excluded from judged answers (cost bias favorable to cheap
teams) and a small truncation audit — are stated but not corrected.

\hrule


\section{7 Future Directions}

Four experiments follow directly. First, a \textbf{subtractive refiner}
study that is permitted to cut as well as add: this isolates the
density mechanism identified in §4.5. Second, a \textit{}human-panel
replication\textit{} on a subset of runs to bound LLM-judge bias in our
specific design. Third, a \textbf{second task domain} (code generation,
scientific writing) to test whether the coverage-vs-density
mechanism generalizes beyond business strategy. Fourth, \textit{}dynamic
pricing integration\textit{}: the qwen off-peak discount and deepseek
peak-surge windows suggest that time-of-day-aware orchestration may
shift the cost curve further; we registered the window but did not
price the study at realized paid prices.

\hrule


\section{8 Conclusion}

For organizations deploying LLMs for knowledge work, three investment
rules emerge from this 1,900-run experiment. \textbf{First}, at the cheap end,
structure substitutes for model tier: a three-step pipeline on a cheap model
delivers 98\% of frontier quality at 23\% of the cost (T-EGR 6.15) — the
highest-return configuration tested. \textbf{Second}, at the frontier end, team
orchestration destroys value: it costs 3.5× the frontier-solo price for
negligible quality gain (T-EGR 0.27). \textbf{Third}, a cheap refinement pass
(\$0.002) recovers most of the frontier gap but never beats it — a bounded-
improvement result that makes it a viable risk-mitigation strategy, not a
replacement strategy. The residual frontier advantage is density (compressed
insight, zero waste), not coverage — a capability gap that add-only pipelines
cannot bridge. Personas are value-destroying: they increase cost while
degrading quality. The governing principle is contingency, not universality:
match the right pipeline architecture to the right model tier for the right
task, and invest in refinement before investing in a frontier team.

\hrule


\section{References}

Albers, C., \& Lakens, D. (2018). When power analyses based on pilot data
are biased: Inaccurate effect size estimators and follow-up bias.
\textit{Journal of Experimental Social Psychology}, 74, 187–195.
DOI: 10.1016/j.jesp.2017.09.004

Brynjolfsson, E., Li, D., \& Raymond, L. R. (2023). \textit{Generative AI at Work}
(NBER Working Paper No. 31161). DOI: 10.3386/w31161

Cohen, J. (1960). A coefficient of agreement for nominal scales.
\textit{Educational and Psychological Measurement}, 20(1), 37–46.
DOI: 10.1177/001316446002000104

Dell'Acqua, F., McFowland III, E., Mollick, E. R., Lifshitz-Assaf, H.,
Kellogg, K., Rajendran, S., Krayer, L., Candelon, F., \& Lakhani, K. R.
(2023). *Navigating the jagged technological frontier: Field experimental
evidence of the effects of AI on knowledge worker productivity and quality*
(HBS Working Paper No. 24-013). DOI: 10.2139/ssrn.4573321

Diewert, W. E. (1976). Exact and superlative index numbers. *Journal of
Econometrics*, 4(2), 115–145. DOI: 10.1016/0304-4076(76)90009-9

Hariri, Y. (2025a). \textit{LLM-Creativity} [dataset and paper]. Zenodo.
DOI: 10.5281/zenodo.17407392

Hariri, Y. (2025b). \textit{LLM Team Composition Paradox} [dataset and paper].
Zenodo. DOI: 10.5281/zenodo.17490945

Hariri, Y. (2026b). *Doing More with Less: How Cultural Personas Bridge the
Cost-Performance Gap in LLMs*. Zenodo. DOI: 10.5281/zenodo.19533619

Lakens, D. (2013). Calculating and reporting effect sizes to facilitate
cumulative science: A practical primer for t-tests and ANOVAs.
\textit{Frontiers in Psychology}, 4, 863.
DOI: 10.3389/fpsyg.2013.00863

Liang, T., He, Z., Jiao, W., Wang, X., Wang, Y., Wang, R., Yang, Y.,
Shi, S., \& Tu, Z. (2024). Encouraging divergent thinking in large language
models through multi-agent debate. In *Proceedings of the 2024 Conference on
Empirical Methods in Natural Language Processing* (pp. 17889–17904).
DOI: 10.18653/v1/2024.emnlp-main.992

Nosek, B. A., Ebersole, C. R., DeHaven, A. C., \& Mellor, D. T. (2018). The
preregistration revolution. \textit{PNAS}, 115(11), 2600–2606.
DOI: 10.1073/pnas.1708274114

Noy, S., \& Zhang, W. (2023). Experimental evidence on the productivity
effects of generative artificial intelligence. \textit{Science}, 381(6654),
187–192. DOI: 10.1126/science.adh2586

Sandve, G. K., Nekrutenko, A., Taylor, J., \& Hovig, E. (2013). Ten simple
rules for reproducible computational research. *PLoS Computational
Biology*, 9(10), e1003285. DOI: 10.1371/journal.pcbi.1003285

Shinn, N., Cassano, F., Berman, E., Gopinath, A., Narasimhan, K., \& Yao, S.
(2023). Reflexion: Language agents with verbal reinforcement learning.
In \textit{Advances in Neural Information Processing Systems 36 (NeurIPS 2023)},
pp. 8634–8652.

Snell, C., Lee, J., Xu, K., \& Kumar, A. (2024). Scaling LLM test-time
compute optimally can be more effective than scaling model parameters.
arXiv:2408.03314. DOI: 10.48550/arXiv.2408.03314

Steegen, S., Tuerlinckx, F., Gelman, A., \& Vanpaemel, W. (2016). Increasing
transparency through a multiverse analysis. *Perspectives on Psychological
Science*, 11(5), 702–712. DOI: 10.1177/1745691616658637

Tomlinson, B., Black, R. W., Patterson, D. J., \& Torrance, A. W. (2024).
The carbon emissions of writing and illustrating are lower for AI than for
humans. \textit{Scientific Reports}, 14, 3732.
DOI: 10.1038/s41598-024-54271-x

Verga, P., Hofstätter, S., Althammer, S., Su, Y., Piktus, A.,
Arkhangorodsky, A., Xu, M., White, N., \& Lewis, P. (2024). Replacing judges
with juries: Evaluating LLM generations with a panel of diverse models.
arXiv:2404.18796. DOI: 10.48550/arXiv.2404.18796

Wang, X., Wei, J., Schuurmans, D., Le, Q., Chi, E., Narang, S.,
Chowdhery, A., \& Zhou, D. (2023). Self-consistency improves chain of thought
reasoning in language models. In *Proceedings of the Eleventh International
Conference on Learning Representations* (ICLR 2023).

Wei, J., Wang, X., Schuurmans, D., Bosma, M., Ichter, B., Xia, F., Chi, E.,
Le, Q., \& Zhou, D. (2022). Chain-of-thought prompting elicits reasoning in
large language models. In *Advances in Neural Information Processing Systems
35 (NeurIPS 2022)*, pp. 24824–24837.

Yao, S., Yu, D., Zhao, J., Shafran, I., Griffiths, T. L., Cao, Y., \&
Narasimhan, K. (2023). Tree of thoughts: Deliberate problem solving with
large language models. In *Advances in Neural Information Processing Systems
36 (NeurIPS 2023)*.

Zheng, L., Chiang, W.-L., Sheng, Y., Zhuang, S., Wu, Z., Zhuang, Y.,
Lin, Z., Li, Z., Li, D., Xing, E. P., Zhang, H., Gonzalez, J. E., \&
Stoica, I. (2023). Judging LLM-as-a-judge with MT-Bench and Chatbot Arena.
In \textit{Advances in Neural Information Processing Systems 36 (NeurIPS 2023)},
Datasets and Benchmarks Track.

\end{document}
